% !TeX program = pdflatex
% =============================================================================
% Vorlage: Praktikum (Laborhandout)
% Kantonsschule KFR – Physik
%
% Verwendung:
%   1. Diese Datei nach praktika/ kopieren und sinnvoll umbenennen,
%      z. B. praktika/freier-fall-praktikum.tex
%   2. Kopfzeile (Klasse/Datum/Thema/Titel) anpassen.
%   3. Lernziele, Materialliste, Theorie, Durchführung und
%      Auswertung/Fragen für das jeweilige Praktikum ausfüllen.
%   4. Lösungen ein-/ausblenden: \printanswers (Lehrerexemplar) bzw.
%      \noprintanswers (Schülerexemplar, Standard) weiter unten umschalten.
%      Erwartete Auswertungsresultate/Musterlösungen gehören in
%      \begin{solution}...\end{solution}, damit sie im Schülerexemplar
%      ausgeblendet bleiben.
% =============================================================================
\documentclass[addpoints,12pt,a4paper]{exam}

\usepackage[T1]{fontenc}
\usepackage[utf8]{inputenc}
\usepackage[ngerman]{babel}
\usepackage{geometry}
\usepackage{amsmath}
\usepackage{siunitx}
\usepackage{tikz}
\usepackage{physics}
\usepackage{circuitikz}
\usepackage{enumitem}

\geometry{a4paper, margin=2.5cm, headheight=14pt}

% --- Lösungen ein-/ausblenden -----------------------------------------------
% Schülerexemplar (Standard): \noprintanswers
% Lehrer-/Musterlösungsexemplar: \printanswers
\noprintanswers
% \printanswers

% --- Punkteformat (falls Auswertungsfragen bewertet werden) -----------------
\pointname{~P.}
\pointsinmargin

% --- Lösungsbox auf Deutsch --------------------------------------------------
\renewcommand{\solutiontitle}{\noindent\textbf{Musterlösung:}\enspace}

% --- Fusszeile auf Deutsch (exam-Klasse zeigt sonst "Page N") ---------------
\cfoot{Seite \thepage}

% --- Kopfzeile: Klasse / Datum / Thema / Titel ------------------------------
% Einheitliches Layout für Arbeitsblatt, Prüfung und Praktikum.
\newcommand{\Kopfzeile}[4]{%
  \begin{center}
    {\Large\bfseries #4}
  \end{center}
  \vspace{0.3em}
  \noindent\textbf{Klasse:}~#1 \hfill \textbf{Datum:}~#2 \hfill \textbf{Thema:}~#3
  \vspace{0.3em}

  \hrule
  \vspace{0.8em}
}

% --- Antwortraum: ca. 2.5 cm Platz pro Punkt für handschriftliche Antworten -
\newlength{\antwortraumlen}
\newcommand{\Antwortraum}[1]{%
  \pgfmathsetlength{\antwortraumlen}{#1*2.5cm}%
  \vspace{\antwortraumlen}%
  \par
}

% --- Praktikums-Abschnittsüberschrift ---------------------------------------
\newcommand{\Abschnitt}[1]{%
  \vspace{0.8em}
  {\large\bfseries #1}
  \vspace{0.3em}

  \hrule
  \vspace{0.5em}
}

% Werte für Kopfzeile zentral setzen
\newcommand{\Klasse}{5b}
\newcommand{\Datum}{XX.XX.20XX}
\newcommand{\Thema}{Mechanik}
\newcommand{\PraktikumsTitel}{Praktikum: Freier Fall}

\begin{document}

\Kopfzeile{\Klasse}{\Datum}{\Thema}{\PraktikumsTitel}

\begin{tabular}{@{}ll@{}}
  \textbf{Name:} \makebox[6cm]{\hrulefill} & \textbf{Gruppe:} \makebox[3cm]{\hrulefill}
\end{tabular}

% --- Lernziele ---------------------------------------------------------------
\Abschnitt{Lernziele}

Nach diesem Praktikum können Sie...
\begin{itemize}[nosep]
  \item die Fallbeschleunigung $g$ experimentell bestimmen.
  \item Messunsicherheiten abschätzen und im Ergebnis angeben.
  \item ein $t$-$s$-Diagramm auswerten, um daraus $g$ zu berechnen.
\end{itemize}

% --- Materialliste ------------------------------------------------------------
\Abschnitt{Materialliste}

\begin{itemize}[nosep]
  \item Fallschnur mit Massestück
  \item Lichtschranken (2×) mit Zeitmessgerät
  \item Massstab / Bandmass
  \item Stativmaterial
\end{itemize}

% --- Theorie ------------------------------------------------------------------
\Abschnitt{Theorie}

Beim freien Fall (ohne Luftwiderstand) gilt für die Fallstrecke $s$ in
Abhängigkeit von der Fallzeit $t$:
\[
  s = \frac{1}{2} g t^2
\]
mit der Fallbeschleunigung $g \approx \SI{9.81}{\meter\per\second\squared}$.
Durch Auftragen von $s$ gegen $t^2$ ergibt sich eine Gerade mit Steigung
$\frac{g}{2}$.

% --- Durchführung ---------------------------------------------------------------
\Abschnitt{Durchführung}

\begin{enumerate}[nosep]
  \item Bauen Sie den Versuchsaufbau gemäss Skizze auf der Wandtafel auf.
  \item Messen Sie die Fallstrecke $s$ zwischen den beiden Lichtschranken
    und notieren Sie sie in Tabelle~1.
  \item Lassen Sie das Massestück fallen und messen Sie die Fallzeit $t$.
    Wiederholen Sie die Messung 5-mal für dieselbe Strecke.
  \item Wiederholen Sie Schritte 2--3 für mindestens 4 weitere Fallstrecken.
\end{enumerate}

\vspace{0.5em}
\noindent\textbf{Tabelle 1: Messwerte}

\begin{center}
\begin{tabular}{|c|c|c|c|c|c|c|}
  \hline
  $s$ [\si{\meter}] & $t_1$ [\si{\second}] & $t_2$ [\si{\second}] & $t_3$ [\si{\second}] & $t_4$ [\si{\second}] & $t_5$ [\si{\second}] & $\overline{t}$ [\si{\second}] \\
  \hline
   & & & & & & \\[1.5em]
  \hline
   & & & & & & \\[1.5em]
  \hline
   & & & & & & \\[1.5em]
  \hline
   & & & & & & \\[1.5em]
  \hline
   & & & & & & \\[1.5em]
  \hline
\end{tabular}
\end{center}

% --- Auswertung / Fragen -------------------------------------------------------
\Abschnitt{Auswertung / Fragen}

\begin{questions}

  \question[3] Tragen Sie $s$ gegen $\overline{t}^2$ in einem Diagramm auf
  (Millimeterpapier oder digital) und bestimmen Sie die Steigung der
  Ausgleichsgeraden.
  \Antwortraum{3}

  \question[3] Berechnen Sie aus der Steigung den Wert für $g$ und
  vergleichen Sie ihn mit dem Literaturwert $\SI{9.81}{\meter\per\second\squared}$.
  \Antwortraum{3}

  \begin{solution}
    Steigung $= g/2$, also $g = 2 \times \text{Steigung}$. Der berechnete
    Wert sollte innerhalb der Messunsicherheit nahe bei
    $\SI{9.81}{\meter\per\second\squared}$ liegen.
  \end{solution}

  \question[2] Nennen Sie zwei mögliche Fehlerquellen in diesem Experiment
  und wie sie sich auf das Ergebnis auswirken.
  \Antwortraum{2}

\end{questions}

\end{document}
